\documentclass[12pt]{article}
\usepackage[brazil]{babel}
\usepackage[utf8]{inputenc}
\usepackage[T1]{fontenc}
\usepackage{amsmath, amssymb}
\usepackage{geometry}
\usepackage{graphicx}
\usepackage{tikz}
\usetikzlibrary{shapes,arrows,positioning,calc,matrix}
\usepackage{pgfplots}
\usepgfplotslibrary{fillbetween}
\pgfplotsset{compat=1.18}
\usepackage{xcolor}
\usepackage{enumitem}
\usepackage{caption}
\usepackage{float}

\definecolor{maincolor}{RGB}{0,100,50}
\usepackage{fancyhdr}
\definecolor{forestgreen}{RGB}{0,100,50}
\usepackage[most]{tcolorbox}

\geometry{margin=2.5cm}
\pagestyle{fancy}
\fancyhf{}

\renewcommand{\headrulewidth}{2pt}
\renewcommand{\footrulewidth}{2pt}

\renewcommand{\headrule}{\hbox to\headwidth{%
  \color{black}\leaders\hrule height \headrulewidth\hfill}}

\renewcommand{\footrule}{\hbox to\headwidth{%
  \color{black}\leaders\hrule height \footrulewidth\hfill}}

\setlength{\headheight}{51pt}

\fancyhead[L]{%
  \includegraphics[height=1.2cm]{Universidade_de_Rio_Verde.jpg}
}

\fancyhead[R]{%
  \raggedleft
  \textcolor{black}{%
    \small
    \textbf{Lista de Exercícios — Cálculo de Complexidade}\\
    Me.~Sergio Souza Novak\\
    \textit{Lógica para Engenharia de Software}
  } 
}

\fancyfoot[C]{%
  \textcolor{black}{\thepage}
}

\newcommand{\quadro}[1]{\vspace{0.2cm}\begin{tcolorbox}[colback=white, colframe=black, height=#1, boxrule=0.5pt, arc=0pt] \textit{Resposta:} \end{tcolorbox}}

\begin{document}

\begin{enumerate}[label=\arabic*), leftmargin=1.5cm, itemsep=0.8cm]

    \item \textbf{Cálculo Linear}: O algoritmo abaixo percorre um vetor $A$ de $n$ elementos para somá-los. Preencha a tabela de custos:
    \begin{center}
        \begin{tikzpicture}
            \draw[thick] (0,0) grid (5,0.6);
            \node at (0.5,0.3) {$A[1]$}; \node at (1.5,0.3) {$A[2]$}; \node at (2.5,0.3) {\dots}; \node at (3.5,0.3) {\dots}; \node at (4.5,0.3) {$A[n]$};
            \draw[->, thick] (0.5,-0.3) -- (0.5,-0.05); \node[below] at (2.5,-0.3) {Percorre as posições sequencialmente};
            \draw[->, thick] (4.5,-0.3) -- (4.5,-0.05);
        \end{tikzpicture}
    \end{center}
    \begin{tcolorbox}[colback=white, colframe=black, fontupper=\small]
        \textbf{SOMA-VETOR($A, n$)} \\
        \begin{tabular}{llcc}
            \hline
            & \textbf{instrução} & \textbf{custo} & \textbf{vezes} \\
            \hline
            1 & $soma = 0$ & \underline{\hspace{1cm}} & \underline{\hspace{1cm}} \\
            2 & \textbf{para} $i = 1$ \textbf{até} $n$: & \underline{\hspace{1cm}} & \underline{\hspace{1cm}} \\
            3 & \quad $soma = soma + A[i]$ & \underline{\hspace{1cm}} & \underline{\hspace{1cm}} \\
            4 & \textbf{retornar} $soma$ & \underline{\hspace{1cm}} & \underline{\hspace{1cm}} \\
            \hline
        \end{tabular}
    \end{tcolorbox}
    \quadro{2cm}

    \item \textbf{Cálculo Quadrático}: Duas matrizes $A$ e $B$ de ordem $n \times n$ são somadas. Analise o impacto dos laços aninhados:
    \begin{center}
        \begin{tikzpicture}[scale=0.8]
            \draw[thick] (0,0) grid (2,2); \node at (1,1) {$A_{n \times n}$};
            \node at (2.5, 1) {$+$};
            \draw[thick] (3,0) grid (5,2); \node at (4,1) {$B_{n \times n}$};
            \node at (5.5, 1) {$=$};
            \draw[thick] (6,0) grid (8,2); \node at (7,1) {$C_{n \times n}$};
            \node[below] at (4, -0.2) {Cada célula $(i,j)$ exige um cálculo};
        \end{tikzpicture}
    \end{center}
    \begin{tcolorbox}[colback=white, colframe=black, fontupper=\small]
        \textbf{SOMA-MATRIZ($A, B, n$)} \\
        \begin{tabular}{llcc}
            \hline
            & \textbf{instrução} & \textbf{custo} & \textbf{vezes} \\
            \hline
            1 & \textbf{para} $i = 1$ \textbf{até} $n$: & \underline{\hspace{1cm}} & \underline{\hspace{1cm}} \\
            2 & \quad \textbf{para} $j = 1$ \textbf{até} $n$: & \underline{\hspace{1cm}} & \underline{\hspace{1cm}} \\
            3 & \quad \quad $C[i,j] = A[i,j] + B[i,j]$ & \underline{\hspace{1cm}} & \underline{\hspace{1cm}} \\
            4 & \textbf{retornar} $C$ & \underline{\hspace{1cm}} & \underline{\hspace{1cm}} \\
            \hline
        \end{tabular}
    \end{tcolorbox}
    \quadro{2cm}


    \item \textbf{Análise Gráfica de Pontos}: Com base no comportamento das curvas no gráfico abaixo, assinale a alternativa que associa corretamente a cor/formato do ponto à sua respectiva complexidade:
    \begin{center}
    \begin{tikzpicture}
        \begin{axis}[
            width=11cm, height=7cm,
            axis lines = left,
            xlabel = {$n$ (Entrada)},
            ylabel = {$T(n)$ (Custo)},
            ymin=0, ymax=1000, xmin=0, xmax=45,
            grid=both,
            grid style={line width=.1pt, draw=black!10},
            legend pos=outer north east,
            legend style={font=\tiny},
            restrict y to domain=0:1100,
        ]
            % n log n - Azul (Círculo)
            \addplot[blue, only marks, mark=*, samples=40, domain=1:40] {x*(ln(x)/ln(2))};
            \addlegendentry{Pontos Azuis ($\bullet$)}
            
            % n^3 - Vermelho (Quadrado)
            \addplot[red, only marks, mark=square*, samples=40, domain=1:40] {x^3};
            \addlegendentry{Pontos Vermelhos ($\blacksquare$)}
            
            % n! - Verde (Triângulo)
            \addplot[green!60!black, only marks, mark=triangle*, samples=40, domain=1:40] {sqrt(2*pi*x)*(x/e)^x};
            \addlegendentry{Pontos Verdes ($\blacktriangle$)}
        \end{axis}
    \end{tikzpicture}
    \end{center}
    
    \begin{enumerate}[label=\alph*)]
        \item Azul: $O(n^3)$; Vermelho: $O(n!)$; Verde: $O(n \log n)$.
        \item Azul: $O(n!)$; Vermelho: $O(n^3)$; Verde: $O(n \log n)$.
        \item Azul: $O(n \log n)$; Vermelho: $O(n^3)$; Verde: $O(n!)$.
        \item Azul: $O(n \log n)$; Vermelho: $O(n!)$; Verde: $O(n^3)$.
    \end{enumerate}

    \item \textbf{Busca em Matriz (Busca Exaustiva)}: Considere uma matriz quadrada $M$ de ordem $n$, que armazena um total de $n^2$ elementos. O algoritmo abaixo realiza uma varredura completa, percorrendo cada linha e cada coluna para localizar um valor específico chamado `alvo`. 
    
    No cenário de \textbf{pior caso}, o elemento buscado está na última célula da última linha (posição $M[n,n]$) ou não está presente na matriz, o que exige que o software verifique absolutamente todas as posições.
    
    \begin{itemize}
        \item \textbf{Exemplo 1}: Em uma matriz $3 \times 3$ ($n=3$), o algoritmo verificará até $9$ elementos ($3^2$).
        \item \textbf{Exemplo 2}: Em uma matriz $100 \times 100$ ($n=100$), o algoritmo verificará até $10.000$ elementos ($100^2$).
    \end{itemize}

    \begin{center}
        \begin{tikzpicture}[scale=0.8]
            \draw[thick] (0,0) grid (3,3);
            \node at (0.5, 2.5) {1}; \node at (2.5, 0.5) {$n^2$};
            \draw[->, thick, dashed] (0.5, 2.5) -- (2.5, 2.5) -- (0.5, 1.5) -- (2.5, 1.5) -- (0.5, 0.5) -- (2.5, 0.5);
            \node[right] at (3.2, 1.5) {Varredura completa};
        \end{tikzpicture}
    \end{center}
    \begin{tcolorbox}[colback=white, colframe=black, fontupper=\small]
        \textbf{BUSCA-MATRIZ($M, n, alvo$)} \\
        \begin{tabular}{llcc}
            \hline
            & \textbf{instrução} & \textbf{custo} & \textbf{vezes} \\
            \hline
            1 & \textbf{para} $i = 1$ \textbf{até} $n$: & \underline{\hspace{1cm}} & \underline{\hspace{1cm}} \\
            2 & \quad \textbf{para} $j = 1$ \textbf{até} $n$: & \underline{\hspace{1cm}} & \underline{\hspace{1cm}} \\
            3 & \quad \quad \textbf{se} $M[i,j] == alvo$: & \underline{\hspace{1cm}} & \underline{\hspace{1cm}} \\
            4 & \quad \quad \quad \textbf{retornar} $(i,j)$ & \underline{\hspace{1cm}} & \underline{\hspace{1cm}} \\
            5 & \textbf{retornar} NULO & \underline{\hspace{1cm}} & \underline{\hspace{1cm}} \\
            \hline
        \end{tabular}
    \end{tcolorbox}
    \quadro{2cm}

    \item \textbf{Laços Independentes vs Aninhados}: O algoritmo abaixo processa o vetor $A$ de duas formas diferentes. Qual termo domina o custo final?

    \begin{tcolorbox}[colback=white, colframe=black, fontupper=\small]
        \textbf{ALGORITMO-MISTO($A, n$)} \\
        \begin{tabular}{llcc}
            \hline
            & \textbf{instrução} & \textbf{custo} & \textbf{vezes} \\
            \hline
            1 & \textbf{para} $i = 1$ \textbf{até} $n$: & \underline{\hspace{1cm}} & \underline{\hspace{1cm}} \\
            2 & \quad print($A[i]$) & \underline{\hspace{1cm}} & \underline{\hspace{1cm}} \\
            3 & \textbf{para} $j = 1$ \textbf{até} $n$: & \underline{\hspace{1cm}} & \underline{\hspace{1cm}} \\
            4 & \quad \textbf{para} $k = 1$ \textbf{até} $n$: & \underline{\hspace{1cm}} & \underline{\hspace{1cm}} \\
            5 & \quad \quad print($A[j] + A[k]$) & \underline{\hspace{1cm}} & \underline{\hspace{1cm}} \\
            \hline
        \end{tabular}
    \end{tcolorbox}
    \quadro{2cm}

    \item \textbf{Infeção de Rede (Ransomware)}: Analise o cenário de um ataque cibernético onde um \textbf{Ransomware} altamente agressivo utiliza uma técnica de replicação recursiva. Ao infectar um servidor, o código malicioso busca vulnerabilidades em outros dois servidores conectados à mesma rede e inicia o processo de infecção neles simultaneamente. Este comportamento se repete para cada novo servidor comprometido, avançando por diversos níveis de profundidade na infraestrutura de TI. 
    
    Qual a complexidade Big-O desta propagação em relação ao \textbf{nível} de profundidade da rede?
    \begin{center}
        \begin{tikzpicture}[scale=0.7]
            \node[draw, circle] (R) at (0,0) {\tiny S1};
            \node[draw, circle] (A1) at (-1.5,-1) {\tiny S2}; \node[draw, circle] (A2) at (1.5,-1) {\tiny S3};
            \node[draw, circle] (B1) at (-2.2,-2) {\tiny S4}; \node[draw, circle] (B2) at (-0.8,-2) {\tiny S5};
            \node[draw, circle] (B3) at (0.8,-2) {\tiny S6}; \node[draw, circle] (B4) at (2.2,-2) {\tiny S7};
            \draw (R)--(A1); \draw (R)--(A2); \draw (A1)--(B1); \draw (A1)--(B2); \draw (A2)--(B3); \draw (A2)--(B4);
        \end{tikzpicture}
    \end{center}
    \begin{tcolorbox}[colback=white, colframe=black, fontupper=\small]
        \textbf{INFECTA-REDE(nível)} \\
        1 \quad \textbf{se} nível $> 0$: \\
        2 \quad \quad Infectar Servidor A (INFECTA-REDE(nível $- 1$)) \\
        3 \quad \quad Infectar Servidor B (INFECTA-REDE(nível $- 1$)) \\
        4 \quad \quad Criptografar arquivos do servidor atual
    \end{tcolorbox}
    \quadro{3cm}
    \item \textbf{Simplificação Assintótica}: Dada a função de custo $T(n) = 5n^2 + 100n + 500$, identifique a complexidade Big-O.
    \quadro{3cm}


    \item \textbf{Busca Binária}: Em um vetor de $n = 1.000.000$ elementos ordenados, a busca binária divide o problema sucessivamente por 2.
    \begin{center}
        \begin{tikzpicture}[scale=0.8]
            \draw[thick] (0,0) rectangle (6,0.5); \node at (3, 0.25) {$n$};
            \draw[thick] (0,-0.8) rectangle (3,-0.3); \node at (1.5, -0.55) {$n/2$};
            \draw[thick] (0,-1.6) rectangle (1.5,-1.1); \node at (0.75, -1.35) {$n/4$};
            \node at (3, -1.35) {\dots};
            \draw[->] (3, 0) -- (1.5, -0.3); \draw[->] (1.5, -0.8) -- (0.75, -1.1);
        \end{tikzpicture}
    \end{center}
    \begin{tcolorbox}[colback=white, colframe=black, fontupper=\small]
        \textbf{BUSCA-BINARIA($A, n, alvo$)} \\
        1 \quad $inicio = 1, fim = n$ \\
        2 \quad \textbf{enquanto} $inicio \leq fim$: \\
        3 \quad \quad $meio = (inicio + fim) / 2$ \\
        4 \quad \quad \textbf{se} $A[meio] == alvo$: \textbf{retornar} $meio$ \\
        5 \quad \quad \textbf{senao se} $A[meio] < alvo$: $inicio = meio + 1$ \\
        6 \quad \quad \textbf{senao}: $fim = meio - 1$ \\
        7 \quad \textbf{retornar} NULO
    \end{tcolorbox}
    Qual a complexidade $O$ e o número máximo de comparações para $n = 1.000.000$?
    \quadro{3cm}

    \item O algoritmo abaixo compara pares adjacentes e os troca. Preencha a tabela de análise:
    \begin{center}
        \begin{tikzpicture}
            \draw[thick] (0,0) grid (4,0.6);
            \draw[<->, thick] (0.5, 0.8) to[bend left] (1.5, 0.8);
            \node at (0.5, 0.3) {8}; \node at (1.5, 0.3) {3}; \node at (2.5, 0.3) {10}; \node at (3.5, 0.3) {1};
            \node[above] at (1, 1.1) {\footnotesize Troca};
        \end{tikzpicture}
    \end{center}
    \begin{tcolorbox}[colback=white, colframe=black, fontupper=\small]
        \textbf{ORDENA($A, n$)} \\
        \begin{tabular}{llcc}
            \hline
            & \textbf{instrução} & \textbf{custo} & \textbf{vezes} \\
            \hline
            1 & \textbf{para} $i = 1$ \textbf{até} $n$: & \underline{\hspace{1cm}} & \underline{\hspace{1cm}} \\
            2 & \quad \textbf{para} $j = 1$ \textbf{até} $n - 1$: & \underline{\hspace{1cm}} & \underline{\hspace{1cm}} \\
            3 & \quad \quad \textbf{se} $A[j] > A[j+1]$: & \underline{\hspace{1cm}} & \underline{\hspace{1cm}} \\
            4 & \quad \quad \quad Troca($A[j], A[j+1]$) & \underline{\hspace{1cm}} & \underline{\hspace{1cm}} \\
            \hline
        \end{tabular}
    \end{tcolorbox}
    \quadro{2cm}


    \item \textbf{Comparação de Curvas}: O algoritmo A possui $100n$ e o algoritmo B possui $n^2$. Determine a partir de qual valor de $n$ o Algoritmo A torna-se mais eficiente.
    \quadro{4cm}

    \item \textbf{Questão de Concurso (CESPE / CEBRASPE - 2025 - Polícia Federal)}: \\
    No que se refere ao SonarQube, às estruturas de dados e à complexidade de algoritmos, julgue o item subsecutivo:
    
    "Para grandes volumes de dados, um algoritmo com complexidade de tempo $O(n)$ (linear) é considerado menos eficiente que um algoritmo com complexidade de tempo $O(n \log n)$, uma vez que o crescimento linear é mais acentuado que o crescimento logarítmico."
    
    \begin{itemize}
        \item[($\quad$)] Certo
        \item[($\quad$)] Errado
    \end{itemize}

    \item \textbf{Questão de Concurso (IV - UFG - 2025 - Câmara de Morrinhos - GO - Analista de TI)}: \\
    Considere o seguinte código escrito em Python 3:
    
    \begin{tcolorbox}[colback=white, colframe=black, fontupper=\small]
\begin{verbatim}
def busca_binaria(lista, elemento):
    inicio = 0
    fim = len(lista) - 1
    
    while inicio <= fim:
        meio = (inicio + fim) // 2
        if lista[meio] == elemento:
            return meio
        elif lista[meio] < elemento:
            inicio = meio + 1
        else:
            fim = meio - 1
            
    return -1
\end{verbatim}
    \end{tcolorbox}
    
    A complexidade de tempo desse algoritmo em termos da notação Big-O é:
    
    \begin{enumerate}[label=\Alph*)]
        \item $O(1)$
        \item $O(n)$
        \item $O(n^2)$
        \item $O(\log n)$
        \item $O(n!)$
    \end{enumerate}

    \item \textbf{Questão de Concurso (COPEVE-UFAL - 2023 - UFAL - Analista de TI)}: \\
    Depois de pensar sobre determinado problema, João fez um rascunho de uma função, produzindo o algoritmo em pseudocódigo do quadro: 
    
    \begin{tcolorbox}[colback=white, colframe=black, fontupper=\small]
        \textbf{função} calculo\_A($n$) \\
        \hspace*{0.5cm} $k = 0$ \\
        \hspace*{0.5cm} \textbf{se} ($n > 1000$) \textbf{então} \\
        \hspace*{1.0cm} \textbf{para} $i$ de 1 até $n$ \textbf{faça} \\
        \hspace*{1.5cm} \textbf{para} $j$ de 1 até $n$ \textbf{faça} \\
        \hspace*{2.0cm} $k = (k * k) + j$ \\
        \hspace*{1.5cm} \textbf{fim para} \\
        \hspace*{1.0cm} \textbf{fim para} \\
        \hspace*{0.5cm} \textbf{senão} \\
        \hspace*{1.0cm} \textbf{para} $j$ de 1 até $n$ \textbf{faça} \\
        \hspace*{1.5cm} $k = (k * j) + 2$ \\
        \hspace*{1.0cm} \textbf{fim para} \\
        \hspace*{0.5cm} \textbf{fim se} \\
        \hspace*{0.5cm} \textbf{retorne} $k$
    \end{tcolorbox}
    
    Qual a complexidade de tempo dessa função na notação Big-O?
    
    \begin{enumerate}[label=\Alph*)]
        \item $O(n)$
        \item $O(n^2)$
        \item $O(n^3)$
        \item $O(\log n)$
        \item $O(n \log n)$
    \end{enumerate}

\end{enumerate}

\end{document}
